% answers_short.tex —— 2013 年·辽宁卷·理科 · 简短答案速查
% 由 tools/make_short_answers.py 自 answers.tex 的 \gksolhead/\gkanswer 提取，
% 并与 2013/2013_数学真题_辽宁理_答案要点.md 逐题比对后生成，勿手改（改源 MD 后重跑）。

\begin{gknotice}
\textbf{使用说明}：本节只给最终答案便于快速核对，完整推导见解析版。答案为编写者独立推导，\textbf{非官方标准答案}。
\end{gknotice}

\gkpart{选择题}
\begin{center}\begin{tabular}{ccccccccccccc}
\toprule
题号 & 1 & 2 & 3 & 4 & 5 & 6 & 7 & 8 & 9 & 10 & 11 & 12 \\
\midrule
答案 & B & D & A & D & B & A & B & A & C & C & B & D \\
\bottomrule
\end{tabular}\end{center}

\gkpart{填空题}
\begin{itemize}[leftmargin=2.6em]
  \item[\textbf{13.}] $16\pi-16$。几何体是底面半径 $2$、高 $4$ 的圆柱挖去底面边长 $2$、高 $4$ 的方柱，体积为 $\pi\cdot2^2\cdot4-2^2\cdot4$。
  \item[\textbf{14.}] $63$。方程两根为 $1$，$4$，递增等比数列取 $a_1=1$，$a_3=4$，公比 $q=2$，故 $S_6=\dfrac{2^6-1}{2-1}=63$。
  \item[\textbf{15.}] $\dfrac57$。由对称性 $|AF|+|BF|=2a$，得 $|BF|=2a-6$；余弦定理定出 $a=7$，$|BF|=8$，$\angle AFB=90^\circ$，故 $c=OF=\dfrac{|AB|}2=5$，$e=\dfrac ca=\dfrac57$。
  \item[\textbf{16.}] $10$。由均值 $7$、方差 $4$ 得总和 $35$、平方和 $265$；柯西不等式给出最大值不超过 $11$，而 $11$ 迫使其余四数全为 $6$，与互不相同矛盾，故最大值为 $10$，由 $4,6,7,8,10$ 实现。
\end{itemize}

\gkpart{解答题}
\begin{itemize}[leftmargin=2.6em]
  \item[\textbf{17.}] (1) $\dfrac{\pi}6$。$|\boldsymbol a|^2=4\sin^2x=1=|\boldsymbol b|^2$，在 $\left[0,\dfrac{\pi}2\right]$ 上得 $x=\dfrac{\pi}6$；(2) $\dfrac32$。$f(x)=\sin\left(2x-\dfrac{\pi}6\right)+\dfrac12$，$x=\dfrac{\pi}3$ 时取最大值 $\dfrac32$
  \item[\textbf{18.}] (1) 证明见解析。由直径所对圆周角为直角得 $BC\perp AC$，又 $PA\perp BC$，故 $BC\perp$ 平面 $PAC$，从而平面 $PBC\perp$ 平面 $PAC$；(2) $\dfrac{\sqrt6}4$。建立空间直角坐标系，两半平面的法向量可取 $\left(\dfrac{\sqrt3}2,\dfrac32,\sqrt3\right)$ 与 $(0,1,0)$，夹角余弦为 $\dfrac{3/2}{\sqrt6}=\dfrac{\sqrt6}4$
  \item[\textbf{19.}] (1) $\dfrac56$。$1-\dfrac{\mathrm C_6^3}{\mathrm C_{10}^3}=1-\dfrac{20}{120}=\dfrac56$；(2) $P(X=0)=\dfrac4{125}$，$P(X=1)=\dfrac{28}{125}$，$P(X=2)=\dfrac{57}{125}$，$P(X=3)=\dfrac{36}{125}$，$E(X)=2$
  \item[\textbf{20.}] (1) $p=2$。切线斜率给出切点参数 $t=-1$，代入 $t^2-2x_0t+4y_0=0$ 并与 $C_2$ 联立，得 $p=2$；(2) $y=\dfrac34x^2$（$x\in\mathbf R$）。设切点参数 $t_1$，$t_2$，则 $x=x_0$，$y=\dfrac{x_0^2-2y_0}2$；用 $y_0=-\dfrac{x_0^2}4$ 消去 $y_0$ 即得
  \item[\textbf{21.}] (1) 证明见解析。令 $h=f-(1-x)$，由 $h''(x)=4x\mathrm e^{-2x}\geqslant0$ 及 $h(0)=h'(0)=0$ 得下界；上界等价于 $(1+x)\mathrm e^{-x}\leqslant1$；(2) $a\leqslant-3$。必要性由 $x=0$ 处 $f-g$ 的右导数 $-(a+3)\geqslant0$ 得出；充分性用 $a\leqslant-3$ 时 $g_a(x)\leqslant g_{-3}(x)$，并验证 $g_{-3}(x)\leqslant1-x\leqslant f(x)$
\end{itemize}

\gkpart{选做题}
\begin{itemize}[leftmargin=2.6em]
  \item[\textbf{22.}] (1) 证明见解析。切线弦定理给 $\angle CEB=\angle EAB$，再由 $AB$ 为直径及 $EF\perp AB$ 得 $\angle FEB=\angle EAB$；(2) 证明见解析。由（1）可证 $\triangle EFB\cong\triangle ECB$，$\triangle EFA\cong\triangle EDA$，故 $AD=AF$，$BC=BF$；再由直角三角形 $AEB$ 的射影定理 $EF^2=AF\cdot BF$
  \item[\textbf{23.}] (1) $(4,\dfrac{\pi}2)$，$(2\sqrt2,\dfrac{\pi}4)$。$C_1$ 为 $x^2+(y-2)^2=4$，$C_2$ 为 $x+y=4$，交点为 $(0,4)$，$(2,2)$；(2) $a=-1$，$b=2$。$P(0,2)$，$Q(1,3)$，直线 $PQ$ 为 $y=x+2$；与参数方程比较斜率和截距即得
  \item[\textbf{24.}] (1) $(-\infty,1]\cup[5,+\infty)$。按 $x<2$，$2\leqslant x\leqslant4$，$x>4$ 三段去掉绝对值求解；(2) $a=3$。不等式化为 $\bigl||x|-|x-a|\bigr|\leqslant1$，因 $a>1$，解集为 $\left[\dfrac{a-1}2,\dfrac{a+1}2\right]$，与 $[1,2]$ 比较得 $a=3$
\end{itemize}
